\chapter{Reliability Qualification Reference}
\label{app:reliability-reference}

This appendix is a lookup reference for qualification standards ownership,
named stress methods, sample and confidence arithmetic, the qualification
planning matrix, and connector and optical-interface test methods. It does not
build the qualification argument. That narrative lives in
\Cref{ch:reliability}: claim, mechanism, stress, observable, acceptance,
samples, confidence, decision. Use this appendix to fill in the method names
and cells once the argument is already framed.

\textit{Read first:} which standard owns which part of the link; what each
named stress method reveals and how it is misread.

\textit{Reference:} sample and confidence rules; planning matrix; connector
methods; Rapid Interview Checks.

\section{Standards ownership in detail}

No single standard covers an optical module. A laser die, a driver IC, a
passive MUX, and an MPO connector are qualified under four different documents
written by four different communities. Split the qualification the same way you
split the failure budget, and state which document backs each claim.

\subsection{Telcordia GR-468 in practice}
\label{sec:gr468}

Optoelectronics inherited a common qualification language from telecom:
\term{Telcordia GR-468-CORE}\citep{gr468}. It covers active optoelectronic
devices, meaning the laser die, the photodiode, and the hermetic or
non-hermetic package around them. The stresses that show up on every laser and
module program are HTOL for powered life or mechanism evidence, temperature
cycling, damp heat, ESD, and mechanical stress. Burn-in appears alongside them
as a production infant-mortality screen when a demonstrated early-life
population justifies it.

Map each stress onto the qualification evidence path in
\Cref{sec:tree-qual-evidence} rather than treating the standard as a second,
parallel sequence. Three practices keep GR-468 evidence honest:

\begin{itemize}
\item Keep burn-in and qualification HTOL distinct. Burn-in screens infant
      mortality from a production population. Qualification HTOL gathers life
      or mechanism evidence from representative samples. A 1{,}000-hour life
      test may justify a room-temperature production proxy, a sampled hot
      audit, a process monitor, or no direct production screen at all. It is
      not itself a per-unit screen.
\item Document the activation energy and confidence bounds whenever HTOL hours
      are converted into field years, and keep sample-size humility
      (\Cref{sec:fit-dppm,ch:manufacturing}).
\item Qualify the laser die, the package, and the module assembly separately
      when failures split across those boundaries
      (\Cref{sec:photonic-packaging,sec:laser-aging,sec:elsfp}).
\end{itemize}

Arrhenius acceleration underpins life projection only when the named failure
mechanism is temperature-accelerated in the assumed regime
(\Cref{sec:accel-validity}).

\subsection{GR-1221: the passive-component companion}

GR-468 covers active optoelectronics. Its companion, \term{Telcordia
GR-1221-CORE} (Generic Reliability Assurance Requirements for Passive Optical
Components), covers the parts GR-468 does not: connectors, fiber couplers, WDM
filters and MUX/DEMUX, splitters, and isolators~\firstcite{gr1221}. It uses the
same style of stress sequence, including damp heat, temperature cycling,
mechanical tests, and aging, but scores pass or fail on insertion loss and
return loss rather than on LIV.

A short-reach link that leans on an on-package or blind-mate MUX and on
external multi-wavelength sources therefore carries a passive reliability
budget that lives in GR-1221, not GR-468
(\Cref{sec:connector-reliability,ch:wdm}).

\subsection{Electronics: JESD47, ESD, latch-up, and AEC-Q100}
\label{sec:ic-reliability}

The modulator driver, TIA, retimer, and DSP (\Cref{sec:drivers,sec:pd-tia}) are
ordinary CMOS or SiGe BiCMOS ICs. They fail by a different and
better-documented set of mechanisms than the laser, so they use the
semiconductor industry's own qualification language rather than laser-aging
math borrowed from \Cref{sec:laser-aging}.

\paragraph{JESD47: the silicon-side GR-468.}
JEDEC \term{JESD47} is the baseline stress-test-driven qualification flow for a
new IC, a device family, or a process change: temperature cycling, HTOL, HTSL,
autoclave or HAST for moisture, and mechanical shock and
vibration~\firstcite{jesd47}. It plays the same role for driver and TIA silicon
that GR-468 plays for the laser. The supplier runs the flow once and the
customer accepts the report instead of renegotiating a qualification plan on
every design win.

\paragraph{ESD and latch-up.}
Two mechanisms are specific to ICs and absent from the optical mechanism
families in \Cref{tab:qual-mechanism-families}:

\begin{description}
\item[ESD] A discharge event during handling or assembly damages a gate oxide
      or junction. Component-level classification uses the human-body model
      (HBM) and charged-device model (CDM) test standards,
      \term{ANSI/ESDA/JEDEC JS-001} and \term{JS-002}~\firstcite{jsesd}. A
      driver or TIA datasheet HBM/CDM rating is the number that protects the
      part on the factory floor, at the fiber-attach and wire-bond stations
      where a laser die is also exposed.
\item[Latch-up] A parasitic thyristor structure in CMOS turns on under an
      overvoltage or current-injection event and holds a low-impedance path
      until power is cycled. \term{JESD78}\firstterm{JESD78} defines the
      overvoltage and $\pm100$~mA current-injection test that classifies
      susceptibility by supply and signal pin~\firstcite{jesd78}. A latched
      driver IC can look like a dead laser on the bench, with no light and no
      LIV signature, until you check the supply current instead of the optical
      path.
\end{description}

ESD and latch-up are primarily design-qualification and process-control
concerns. Targeted production measurements or screens may be used when their
detection capability is validated, but many latent ESD conditions cannot be
economically detected on every finished unit. Do not project either mechanism
with an activation energy. A field ESD or latch-up failure is a manufacturing
or design-margin item (\Cref{sec:fleet-triage}), not a wear-out FIT argument.

\paragraph{AEC-Q100 is supplemental evidence, not a datacenter requirement.}
\term{AEC-Q100}\firstterm{AECQ100} is the automotive industry's qualification
standard for ICs, built on the same JEDEC JESD47 and JESD22 stress methods with
tighter ESD targets and named temperature grades from Grade~3 ($-40$ to
$85$\textdegree C) up to Grade~0 ($-40$ to $150$\textdegree
C)~\firstcite{aecq100}. Datacenter optics does not require Q100. The fleet
lives in a controlled data hall, not an engine bay, and no datacenter
transceiver specification calls for automotive grading.

Treat a published Q100 grade as a useful supplemental signal instead. A driver,
TIA, or retimer die that also ships in an automotive part number carries a
grade that is a fast proxy for the ESD and latch-up margin and the
temperature-cycle depth behind the datasheet, which saves you from re-running
the supplier's qualification plan yourself. It does not replace a product-level
mechanism argument, and its absence is not a finding.

\paragraph{Where IC qualification lands after release.}
Carry IC-level qualification into the production acceptance and SPC structure
in \Cref{ch:manufacturing}. Require the supplier's JESD47 report and the
HBM/CDM/latch-up ratings for driver and TIA die at DVT. Add an ESD handling
audit to the incoming-QC checklist alongside laser LIV and SMSR sampling. Treat
a driver or TIA silicon revision the same way you treat a laser die revision or
a CMIS firmware revision: an ECO that needs first-article requalification, not
a silent BOM swap.

\section{Stress-method quick reference}
\label{sec:stress-method-reference}

Each method below answers a specific question. None of them is evidence on its
own. The evidence is the pairing of a named mechanism with a stress that
accelerates it, an observable that reveals the change, and an acceptance limit
set before the stress starts.

Durations, temperatures, humidity levels, cycle counts, and acceleration
factors are product-specific and are set by the claim, the standard in force,
and the customer requirement. Do not carry a number from one program into
another without re-deriving it.

\begin{table*}[ht]
\refstepcounter{table}\label{tab:stress-method-reference}%
\centering
\setlength{\tabcolsep}{3pt}%
\footnotesize
\renewcommand{\arraystretch}{1.2}%
\begin{tabular}{@{}p{0.15\linewidth}p{0.18\linewidth}p{0.07\linewidth}p{0.28\linewidth}p{0.28\linewidth}@{}}
\toprule
Method & Mechanism family & Powered & Typical observable & Common misuse \\
\midrule
HTOL & Powered semiconductor degradation & Yes & $I_\mathrm{th}$, slope, $\lambda$, OMA, BER, bias creep & Treated as a per-unit ship screen, or as burn-in \\
\addlinespace[0.6em]
HTSL (storage life) & Unbiased thermal and material change & No & Post-bake parametric shift, bond and epoxy integrity & Read as a powered-life result \\
\addlinespace[0.6em]
Temperature cycling & Thermo-mechanical fatigue & Either & Continuity, intermittent lanes, coupling shift, BER & Final functional pass only, so transient opens are missed \\
\addlinespace[0.6em]
Thermal shock & Fast-ramp thermo-mechanical fatigue & Usually no & Cracks, delamination, alignment step & Substituted for cycling when the ramp rate is not the threat \\
\addlinespace[0.6em]
Damp heat & Moisture, corrosion, material degradation & Either & Leakage, insertion loss, ORL, seal integrity & Described as a waterproofing test \\
\addlinespace[0.6em]
HAST & Accelerated moisture ingress & Either & Same as damp heat, at shorter exposure & Pressure and temperature drive a mechanism the field never sees \\
\addlinespace[0.6em]
Mechanical shock & Discrete acceleration events & Usually no & Alignment step, opens, physical damage & Assumed to cover repeated excitation \\
\addlinespace[0.6em]
Vibration & Repeated mechanical excitation & Often in situ & Intermittent contacts, lane drops, loss modulation & Run without in-situ monitoring \\
\addlinespace[0.6em]
Connector mate cycling & Optical-interface wear and debris & No & Insertion loss growth, ORL, endface grade & Datasheet cycle rating inherited without service-rate check \\
\addlinespace[0.6em]
ESD (HBM/CDM) & Electrical overstress & No & Pass/fail classification level & Projected with an activation energy \\
\addlinespace[0.6em]
Latch-up (JESD78) & Parasitic turn-on under injection & Yes & Supply-current latch, recovery on power cycle & Diagnosed as a dead laser \\
\bottomrule
\end{tabular}
\par\medskip
{\raggedright\footnotesize\textbf{Table~\thetable.} Named stress methods and
what each one reveals. Durations, levels, and cycle counts are product-specific
and are not implied here. Acceptance is a bounded, pre-declared change tied to
the life claim, not merely a post-stress functional pass. Standards ownership:
\Cref{sec:gr468,sec:ic-reliability}. Argument structure:
\Cref{sec:qual-argument}.\par}
\end{table*}
\FloatBarrier

\section{Sample size, confidence, and evidence sufficiency}
\label{sec:qual-sample-reference}

Sample strategy is set by the failure-rate or degradation target, the required
confidence, the cost of the units, and the variation in the population. The
narrative rules are in \Cref{sec:fit-dppm,sec:qual-sample-strategy}; the
mechanics are below.

\begin{description}
\item[Zero-failure upper bound] Zero failures establishes a one-sided upper
      confidence bound on the failure rate under the stated assumptions, not a
      zero rate. State the bound and the confidence level, or do not make the
      claim.
\item[Sample-hours] Evidence scales with units multiplied by stress hours,
      adjusted by the acceleration factor. Twenty units for 1{,}000 hours and
      1{,}000 units for 20 hours are not interchangeable, because they expose
      different amounts of population variation.
\item[Censoring] A test stopped at a fixed time, or units removed before the
      end, produces censored data. Say which units ran to what exposure and how
      the censoring was handled before quoting a rate.
\item[Lot and site diversity] Cover lots, date codes, suppliers, manufacturing
      sites or lines, and process corners. A single-lot result bounds that lot.
      It does not bound the population the factory will ship.
\item[Degradation versus failure] Parametric drift is usable evidence long
      before a hard failure appears, and it needs far fewer units. Record
      intermediate reads and fit the trend rather than waiting for units to
      die.
\item[Acceleration factor applicability] An acceleration factor applies to one
      mechanism in one regime. Do not apply a single factor across laser wear,
      corrosion, solder fatigue, and unrelated mechanisms
      (\Cref{sec:accel-validity}).
\end{description}

\heuristic{Write the sentence you intend to claim before you size the sample.
If the sentence is ``no more than $X$ failures per $10^{9}$ device-hours at
$Y$\% confidence, for this population, over this exposure,'' the sample size
follows from arithmetic. If the sentence is ``the product is reliable,'' no
sample size will support it.}

\section{Qualification planning matrix}
\label{sec:qual-planning-matrix}

The matrix restates the worked argument in \Cref{sec:qual-worked-laser} as
mechanism, stress, observable, acceptance, and production control, then repeats
the pattern for the other mechanism families. Fill the cells for the product
class and claimed life. A blank row is an uncovered mechanism, not a covered
one.

\begin{table*}[ht]
\refstepcounter{table}\label{tab:qual-planning-matrix}%
\centering
\setlength{\tabcolsep}{3pt}%
\footnotesize
\renewcommand{\arraystretch}{1.15}%
\begin{tabular}{@{}p{0.18\linewidth}p{0.18\linewidth}p{0.18\linewidth}p{0.18\linewidth}p{0.20\linewidth}@{}}
\toprule
Failure mechanism & Stress & Observable & Acceptance & Production control \\
\midrule
Laser degradation & Temperature and bias (HTOL) & $I_\mathrm{th}$, slope, $\lambda$, BER & Named drift limit vs life claim & ATP proxy, SPC, sampled audit, or none \\
\addlinespace[0.7em]
Solder and attach fatigue & Temperature cycling & Resistance, BER, opens & Post-stress continuity and BER & Process control, FAIR \\
\addlinespace[0.7em]
Contamination and corrosion & Damp heat or HAST & Loss, ORL, leakage & IL/ORL and functional limits & Handling, sealing, audit \\
\addlinespace[0.7em]
Connector wear & Mate cycling & Insertion loss, ORL, endface grade & Cycle-count IL and ORL budget & Supplier rating, service hygiene \\
\addlinespace[0.7em]
Driver, TIA, and DSP silicon & JESD47 flow; HBM/CDM; JESD78 & Parametric shift, supply current, lane failure & Supplier qual report plus product-level limits & Incoming QC, ESD handling audit, ECO control \\
\addlinespace[0.7em]
Package and passive optics & GR-1221-style sequence; shock and vibration & IL, ORL, alignment step, seal integrity & Passive loss budget over life & Assembly SPC, first-article, sampled DPA \\
\bottomrule
\end{tabular}
\par\medskip
{\raggedright\footnotesize\textbf{Table~\thetable.} Qualification planning
matrix. This is a reference aid for cell filling and coverage review. It does
not replace the mechanism argument in \Cref{sec:qual-argument}, and a filled
row is only evidence when the stress is justified for that mechanism.
Interview form: \Cref{fw:qual-plan,sec:tree-qual-evidence}.\par}
\end{table*}
\FloatBarrier

\section{Connector and optical-interface reference}
\label{sec:connector-reliability}

Multi-fiber connectors are the highest-touch mechanical interface in the fleet.
Every ELSFP swap, every fiber-array rework, and every cable-plant install mates
and unmates an MPO\firstterm{MPO}. The MPO/MT ferrule family (rectangular, 6.4~mm $\times$
2.5~mm, guide-pin aligned, 8, 12, 16, or 24 fibers per row) is standardized in
\term{IEC 61754-7}, split into one-fibre-row and two-fibre-row
parts~\firstcite{iec617547}.

That standard fixes geometry, not lifetime. Lifetime comes from two companion
test methods. \term{IEC 61300-2-2} specifies the mate and unmate cycling test
that connector datasheets are rated against, and \term{IEC 61300-3-35} grades
endface scratches, pits, and debris into pass/fail zones on the fiber core and
cladding~\firstcite{iec61300}. TIA-568.3 sets 500 cycles as the
structured-cabling mating-durability floor. MPO and MTP-class connectors are
commonly rated well above 1000 cycles in practice, but that headroom erodes
quickly under poor cleaning discipline (\Cref{sec:optical-channel}).

Three practical consequences follow for an ELSFP or CPO fiber-attach program.

\begin{enumerate}
\item ORL creep is a mating-cycle and cleaning problem before it is a laser
      problem. A rising RIN floor after repeated ELS swaps
      (\Cref{tab:qual-mechanism-families}) is diagnosed with an
      IEC~61300-3-35-style endface inspection, not a laser FA request.
\item Mate-cycle count belongs in the same telemetry you already read for CMIS
      and DDM (\Cref{sec:cmis}). Track it per connector rather than per module,
      since a connector can outlive several module swaps or the reverse.
\item Acceptance limits may include an explicit mating-cycle count and endface
      grade rather than an inherited generic MPO datasheet number
      (\Cref{tab:atp-laser}). An ELS bank that hot-swaps weekly reaches a
      500-cycle floor in under ten years, and a CPO fiber array that is
      field-serviced more aggressively reaches it faster.
\end{enumerate}

ELSFP cycling adds connector wear and contamination that raise ORL
(\Cref{sec:optical-channel,sec:elsfp}). The mating-cycle and endface-grade
numbers above are what turn ``the connector feels loose'' into a measurable
limit instead of a guess.

Destructive physical analysis (cross-section, EDX) and structured 8D or CAPA
with suppliers close the loop from RMA to design rule
(\Cref{sec:supplier-exec,sec:fleet-triage}). Without that loop, packaging
failures get mis-attributed to laser Arrhenius models and the wrong part gets
redesigned.

\section{Rapid Interview Checks}

\paragraph{Prompt.} Why is HTOL not burn-in?\\
\textit{Check.} HTOL gathers life or mechanism evidence from representative
samples. Burn-in screens a production population for a demonstrated early-life
mechanism. Different population, different question, different decision
(\Cref{sec:qual-terms}).

\paragraph{Prompt.} What does zero failures establish?\\
\textit{Check.} A one-sided upper confidence bound on the failure rate for the
tested population, exposure, and assumptions. Not a zero rate, and not a fleet
claim (\Cref{sec:fit-dppm}).

\paragraph{Prompt.} When can Arrhenius be used?\\
\textit{Check.} When the named mechanism is temperature-activated in the
assumed regime, the activation energy comes from that mechanism, and the stress
did not introduce a different failure physics (\Cref{sec:accel-validity}).

\paragraph{Prompt.} What does a standards pass fail to establish?\\
\textit{Check.} That the mechanisms relevant to this product were covered, that
the samples represented the shipping population, and that the remaining margin
supports the life claim. A standard supplies methods and shared language, not
the argument.

\paragraph{Prompt.} Why can a post-stress functional pass miss fatigue?\\
\textit{Check.} Temperature cycling and vibration can open a joint or shift an
alignment transiently. The unit recovers at room temperature, so only in-situ
or periodic monitoring catches it (\Cref{tab:stress-method-reference}).

\paragraph{Prompt.} Why is stable average power insufficient after connector
cycling?\\
\textit{Check.} Reflections raise the RIN and BER floor while average power
barely moves. Inspect the endface and measure ORL before clearing the connector
(\Cref{sec:connector-reliability}).
